\documentclass[12pt]{article}
\usepackage{amsmath,amssymb}

\begin{document}
\section*{{2025 AIME II Problems/Problem 15}}
Source: AoPS Wiki

\subsection*{Solution}
Let $n$ be the minimum value of the expression (changes based on the value of $k$ , however is a constant). Therefore, we can say that
\begin{align*}
f(x)-n=\frac{(x-\alpha)^2(x-\beta)^2}{x}
\end{align*}
This can be done because $n$ is a constant, and for the equation to be true for all $x$ , the right side is also a quartic. The roots must also both be double, or else there is an even smaller value, a contradiction. We expand as follows, comparing coefficients: \begin{align*}
(x-18)(x-72)(x-98)(x-k)-nx=(x-\alpha)^2(x-\beta)^2 \\
-2\alpha-2\beta=-18-72-98-k \implies \alpha+\beta=94+\frac{k}{2} \\
\alpha^2+4\alpha \beta +\beta^2=(-18\cdot -72)+(-18\cdot-98)+(-18\cdot-k)+(-72\cdot-98)+(-72\cdot-k)+(-98\cdot-k)=10116+188k \\
(\alpha^2)(\beta^2)=(-18)(-72)(-98)(-k) \implies \alpha \beta=252\sqrt{2k} \\
\end{align*} Recall $(\alpha+\beta)^2+2\alpha \beta=\alpha^2+4\alpha \beta +\beta^2$ , so we can equate and evaluate as follows: \begin{align}
(94+\frac{k}{2})^2+504\sqrt{2k}=10116+188k \tag{1}\\
\end{align}
\begin{align*}
(47-\frac{k}{4})^2+126\sqrt{2k}=2529 \\
\frac{k^2}{16}-\frac{47}{2}k+126\sqrt{2k}-320=0 \\
\end{align*} We now have a quartic with respect to $\sqrt{k}$ . Keeping in mind it is much easier to guess the roots of a polynomial with integer coefficients, we set $a=\frac{k}{8}$ . Now our equation becomes \begin{align*}
4a^2-188a+504\sqrt{a}-320=0 \\
a^2-47a+126\sqrt{a}-80=0 \\
\end{align*} We can test for easy roots and find $\sqrt{a}=1$ and $2$ . After this, solving the resulting quadratic gets you the remaining roots as $5$ and $-8$ . Discard $-8$ . Working back through our substitution for $a$ , we have generated values of $k$ as $(8, 32, 200)$ . The sum of all $k$ then must be $8+32+200=\boxed{240}$ . ~ lisztepos ~ Edited by aoum Note:
I'm not sure what the first author meant by, "or else there is an even more 'minimum' value." The most noticeable reason there are two double roots is that there are two distinct positive solutions per the conditions of the problem ~ fermat_sLastAMC Further notes (from the author):
To clarify why the equation takes the form above, if you do not regard the context of the problem, you could say that $\frac{(x-\alpha)^3(x-\beta)}{x}$ is another form, also having two roots. However when we consider this form, (WLOG assuming $\alpha < \beta$ ) if we take a value such that $\alpha < x < \beta$ , then $x-\beta < 0$ , therefore implying that $\frac{(x-\alpha)^3(x-\beta)}{x}<0$ or otherwise resulting $f(x)<n$ , which is contradiction as we assumed $n$ is the least value of the function. Similarly, logic proceeds if the third power is on the other binomial.

\end{document}
